\documentclass[10pt,letterpaper]{article}
\usepackage{cornell-notes}

\title{Chapter 20: Chromatic Polynomial}
\author{}
\date{}

\setDocTitle{Chapter 20: Chromatic Polynomial}
\setDocSubtitle{Combinatorial Algorithms}
\setDocVersion{v1.0}
\setDocStatus{Study Companion}
\setDocOwner{Combinatorial Algorithms Cornell Notes}
\setDocAuthor{}
\setDocDate{\today}
\setDocSummary{Cornell-format study and review notes for Chapter 20: Chromatic Polynomial.}

\setCornellCollection{Combinatorial Algorithms Cornell Notes}
\setCornellUnitType{Chapter}
\setCornellUnitNumber{20}
\setCornellUnitTitle{Chromatic Polynomial}
\setCornellCompanionLabel{Study and review companion}

\hypersetup{pdftitle={Chapter 20: Chromatic Polynomial},pdfsubject={Cornell notes},pdfauthor={}}

\begin{document}
\maketitle
\begin{tcolorbox}[colback=Paper,colframe=Ink]{\small\bfseries\color{Teal}CORNELL NOTES}\par{\LARGE\bfseries Chapter 20: Chromatic Polynomial of a Graph (CHROMP)}\par\medskip\textbf{Topic:} Deletion-identification search\hfill\textbf{Date:} \today\par\textbf{Study purpose:} Reduce a connected graph to tree leaves and accumulate polynomial coefficients.\end{tcolorbox}
\begin{tcolorbox}[colback=Teal!5,colframe=Teal,title=Essential question]How does deletion-identification turn proper-coloring counts into a binary graph search ending at easily scored trees?\end{tcolorbox}
\begin{tcolorbox}[colback=white,colframe=Mist,title=Learning objectives]\begin{itemize}\item Define the chromatic polynomial.\item Prove the deletion-identification recurrence.\item Explain why nonforest edges are safe deletion choices.\item Accumulate the tree-basis coefficients.\item Convert those coefficients with POLY and state the exponential cost.\end{itemize}\textbf{Prerequisites:} proper colorings, connected graphs, spanning trees, stack-based search.\end{tcolorbox}
\section*{Cornell notes}
\begin{longtable}{@{}Q!{\color{Mist}\vrule width 1pt}A@{}}\cornellhead
\cue{What does $P(\lambda;G)$ count?}{For a nonnegative integer $\lambda$, it counts proper vertex colorings of $G$ with $\lambda$ available colors. These values agree with a degree-$n$ polynomial when $G$ has $n$ vertices.}
\cue{What is the reduction formula?}{For edge $uv$, let $G-uv$ delete the edge and let $G_{uv}$ identify its endpoints. Colorings of $G-uv$ either color $u,v$ differently, producing a coloring of $G$, or equally, corresponding to a coloring of $G_{uv}$. Hence
\[
P(\lambda;G)=P(\lambda;G-uv)-P(\lambda;G_{uv}).
\]}
\cue{Why stop at trees?}{Every tree on $m$ vertices has
\[
t_m(\lambda)=\lambda(\lambda-1)^{m-1},
\]
because the root has $\lambda$ choices and every other vertex has $\lambda-1$ choices after its parent. Thus each terminal graph can be scored by its vertex count alone.}
\cue{Which edge is selected?}{Call SPANFO so a spanning tree occupies the first $m-1$ edge positions. Any later edge lies outside the tree; deleting it cannot disconnect the graph. Repeating deletion or endpoint identification eventually reaches a tree on every search branch.}
\cue{How is the search organized?}{For a non-tree graph, form the identified child for immediate processing and push the deletion child for later work. When a tree is reached, increment $b_m$, where $m$ is its number of vertices, then pop the next saved graph. This depth-first stack explores every I/D branch.}
\cue{What form is accumulated?}{If $b_m$ is the signed net number of terminal trees of order $m$ produced by the recurrence, then
\[
P(\lambda;G)=\sum_{m=1}^{n} b_m\lambda(\lambda-1)^{m-1}.
\]
The chapter notes the resulting nonnegative coefficients in this tree basis.}
\cue{How are other coefficients obtained?}{POLY converts the tree-basis expression to ordinary power coefficients and to a factorial basis. The latter relates coefficients to the number of colorings using exactly a specified number of colors.}
\cue{Cost and storage}{The search has at most an exponential number of graph states in the edge count. With $O(E)$ work per state, the stated bound is $O(E2^E)$. A stack stores pending graphs with decreasing vertex and edge counts.}
\cue{Cautions}{Input is a connected graph and the endpoint list is destroyed. Identification must remove duplicate edges created by merging endpoints. The exponential bound makes the routine appropriate only for modest graphs.}
\cue{Retrieval practice}{\begin{enumerate}\item What two classes partition colorings of $G-uv$?\item Why is a non-tree edge safe to delete?\item What is a tree's chromatic polynomial?\item What does $b_m$ count?\item Why is the algorithm exponential?\end{enumerate}}
\end{longtable}
\begin{tcolorbox}[colback=Ink!4,colframe=Ink,title=Cornell summary]
CHROMP repeatedly applies the deletion-identification identity for proper colorings. A spanning tree identifies an edge whose deletion preserves connectivity; the algorithm explores both deleting that edge and identifying its endpoints. Every branch eventually reaches a tree, whose chromatic polynomial is the known basis function $\lambda(\lambda-1)^{m-1}$. A depth-first stack records deletion branches while identification branches continue immediately, and terminal trees accumulate coefficients by vertex count. POLY can then transform the tree-basis result into ordinary or factorial coefficients. The method is constructive and exact, but its graph-state tree is exponential in the number of edges and it destructively rewrites endpoint lists.
\end{tcolorbox}
\end{document}
